% Source-derived chapter 05: Variations of Hodge structure and the Kodaira-Parshin trick
% Original source: geometric-local-systems-general-curves-isomonodromy
\chapter{Variations of Hodge structure and the Kodaira-Parshin trick}
\label{geometric-local-systems-general-curves-isomonodromy-chapter-05}
\source{geometric-local-systems-general-curves-isomonodromy}

\begin{remark}[Source section]
\label{geometric-local-systems-general-curves-isomonodromy-chapter-05-source}
\source{geometric-local-systems-general-curves-isomonodromy}
\source{geometric-local-systems-general-curves-isomonodromy}
This chapter reproduces the corresponding section of the retrieved
LaTeX source, including its definitions, statements, examples, and
proofs. It has not yet been translated into Lean.
\notready
\end{remark}

% The source section title was: Variations of Hodge structure and the Kodaira-Parshin trick
\label{section:counterexample}
\source{geometric-local-systems-general-curves-isomonodromy}

In this section we find that
variations of Hodge structure on $\mathscr{M}_{g,1}$ with monodromy which is
``big" in a suitable sense provide examples of flat vector bundles on curves
whose isomonodromic deformation to a nearby curve is never semistable. We then produce such variations of
Hodge structure via the Kodaira-Parshin trick. 
This will be used to prove
\autoref{theorem:counterexample} and contradicts earlier claimed theorems of
Biswas, Heu, and Hurtubise 
\cite[Theorem 1.3]{BHH:logarithmic},
\cite[Theorem 1.3]{BHH:irregular}, and
\cite[Theorem 1.2]{BHH:parabolic},
as described further in \autoref{remark:bhh-error}.

In \autoref{section:hodge-theoretic-results}, we will use that
variations of Hodge structure with suitably large monodromy yield flat vector
bundles which do not have isomonodromic deformations to semistable bundles.
This will be used to analyze variations of Hodge structure on an analytically very general curve.
\subsection{Construction of the counterexample}
We now set up the proof of \autoref{theorem:counterexample} and
\autoref{corollary:counterexample}. We will construct a variation of Hodge
structure over the analytic moduli stack $\mathscr{M}_{g,1}$ whose restriction to each fiber of the forgetful map $\mathscr{M}_{g,1}\to \mathscr{M}_g$ satisfies the hypotheses of \autoref{lemma:positive-Hodge-bundle}. We will do this via the Kodaira-Parshin trick
(see
\cite[Proposition 7]{parshin:algebraic-curves-over-function-fields-i}
and also
\cite[Proposition 7.1]{lawrenceV:diophantine-problems}), 
which produces a family of curves over $\mathscr{M}_{g,1}$ which is non-isotrivial when restricted to each fiber of the natural forgetful map $\mathscr{M}_{g,1}\to \mathscr{M}_g$. 
We give a proof appealing to \cite{cataneseD:answer-to-a-question-by-fujita},
but one can also prove it using \autoref{prop:basic-facts}  and
\autoref{corollary:unstable}, as we mention in
\autoref{remark:alternate-hodge-proof}
Because we had difficulty finding a suitable reference, we now present a version
of the Kodaira-Parshin trick in families.

\begin{proposition}[Kodaira-Parshin trick]
\source{geometric-local-systems-general-curves-isomonodromy}
\notready
	\label{proposition:kodaira-parshin}
\source{geometric-local-systems-general-curves-isomonodromy}
	Let $Y$ denote a Riemann surface of genus $g\geq 1$ with a point $p \in Y$ and
	let $Y^{\circ} := Y - p$. Choose a basepoint $y \in Y^{\circ}$.
	Suppose $G$ is a finite center-free group with a surjection
	$\phi: \pi_1(Y^{\circ},y) \twoheadrightarrow G$ which sends the loop around the puncture $p \in Y$ to 
	a non-identity element of $G$.
	Then there exists a smooth proper relative dimension $1$ map of
	analytic stacks $f: \mathscr{X} \to
	\mathscr{M}_{g,1}$ so that the fiber over a geometric point $[(C,p)] \in
	\mathscr{M}_{g,1}$ is a finite disjoint union of $G$-covers of $C$
	ramified at $p$.
\end{proposition}
We will prove this below in \autoref{subsubsection:proof-kodaira-parshin}.
\begin{remark}
\source{geometric-local-systems-general-curves-isomonodromy}
\notready
	\label{remark:}
\source{geometric-local-systems-general-curves-isomonodromy}
	In the finite disjoint union of $G$-covers appearing at the end of the
	statement of \autoref{proposition:kodaira-parshin},
	we can explicitly identify the finite set of $G$-covers.
	Namely, suppose $h \in \pi_1(\mathscr M_{g,2})$, viewed as an
	automorphism of the fundamental group of a $2$-pointed genus $g$ curve $\pi_1(Y^{\circ},y)$. (See \autoref{subsubsection:setup} below for an explanation of the action of $\pi_1(\mathscr{M}_{g,2})$ on $\pi_1(Y^\circ, y)$.)
	There is one cover associated to each map
	of the form $\phi_h: \pi_1(Y^{\circ},y) \to G$, with $\phi_h(g) :=
	\phi(hgh^{-1})$,
	modulo the following equivalence relation:
	we identify $\phi_h \sim \phi_{g}$ if they
	are conjugate, i.e., if there exists $m \in G$ with $\phi_h(s) = m
	\phi_g(s)m^{-1}$ for all $s \in \pi_1(Y^{\circ},y)$.
\end{remark}


\subsubsection{Setup to prove \autoref{proposition:kodaira-parshin}}
\label{subsubsection:setup}
\source{geometric-local-systems-general-curves-isomonodromy}
Let $$\pi:\mathscr{M}_{g, 2}\to \mathscr{M}_{g,1}$$ be the natural forgetful map.
Let $x\in \mathscr{M}_{g,2}$ be a point. Let $\bar{x} := \pi(x) \in
\mathscr{M}_{g,1}$ and let $C^\circ\subset \mathscr{M}_{g,2}$
denote the fiber of $\pi^{-1}(\bar{x})$.

Note that $C^\circ$ is the complement of a point in a smooth proper connected curve of genus $g$. 
There is a natural short exact sequence
\begin{align}
1\to \pi_1(C^\circ, x) \to \pi_1(\mathscr M_{g,2}, x) \to \pi_1(\mathscr
M_{g,1}, \bar x)\to 1
\end{align}
associated to the map $\mathscr{M}_{g,2} \to \mathscr{M}_{g,1}$ with fiber
$C^{\circ}$.
We may obtain this from the Birman exact sequence \cite[Theorem
4.6]{farbM:a-primer} for mapping class groups,
after identifying the fundamental group for $\mathscr M_{g,n}$ with the mapping
class group of an $n$-times punctured genus $g$ surface.
(The case $n = 0$ follows from contractibility of the universal cover of
$\mathscr M_g$ \cite[Theorem 10.6]{farbM:a-primer} with covering group given by
the mapping class group,
and the case of general $n$ can be deduced from the Birman exact sequence
\cite[Theorem 4.6]{farbM:a-primer}.)

Let $G$ be a center-free finite group and suppose further there
is a surjection $$\gamma: \pi_1(C^\circ, x)\twoheadrightarrow G.$$
We assume that $\gamma$ takes the conjugacy class of the loop around the puncture of $C^\circ$ to a non-identity conjugacy class of $G$.

Define $\Gamma\subset \pi_1(\mathscr{M}_{g,2}, x)$ as the set of $h \in
\pi_1(\mathscr{M}_{g,2},x)$ such that there exists $\widetilde{\gamma}(h)\in G$ with $$\gamma(hgh^{-1})=\widetilde{\gamma}(h)\gamma(g)\widetilde{\gamma}(h)^{-1}$$ for all $g\in \pi_1(C^\circ, x).$ 

\begin{lemma}
\source{geometric-local-systems-general-curves-isomonodromy}
\notready
	\label{lemma:center-free-group}
\source{geometric-local-systems-general-curves-isomonodromy}
	Keeping notation as in \autoref{subsubsection:setup},
	the map $\gamma$ determines a well-defined surjective homomorphism 
	\begin{align*}
		\widetilde{\gamma}: \Gamma  & \rightarrow G\\
		h & \mapsto \widetilde{\gamma}(h).
	\end{align*}
	Moreover, $\Gamma \subset\pi_1(\mathscr{M}_{g,2},x)$ has finite index.
\end{lemma}
\begin{proof}
	We first claim that $\Gamma$ contains $\pi_1(C^\circ, x)$ and surjects
	onto $G$. Indeed, for $h\in
\pi_1(C^\circ, x)$, one may take $\widetilde{\gamma}(h)=\gamma(h)$.
Therefore, the surjectivity of $\gamma$ implies that $\widetilde{\gamma}$ is
also surjective.

	Next, we claim that for each $h$, $\widetilde{\gamma}(h)$ is uniquely determined. 
	Indeed, suppose $\widetilde{\gamma}(h)$ may be either
	$\alpha$ and $\beta$. Then we would have $\alpha \gamma(g) \alpha^{-1} =
	\beta \gamma(g) \beta^{-1}$.
	Since $\gamma$ is surjective, as shown above, we find $\alpha
	\beta^{-1}$ lies in the center of $G$, and therefore is trivial. So
	$\alpha = \beta$.	

	The uniqueness of $\widetilde{\gamma}(h)$ just established shows that $\widetilde{\gamma}$ determines a well-defined map.
This is moreover a homomorphism by the above established uniqueness, because we
then obtain $\widetilde{\gamma}(h) \widetilde{\gamma}(h') =
\widetilde{\gamma}(hh')$.   

Finally, we claim $\Gamma$ has finite index in $\pi_1(\mathscr{M}_{g,2}, x)$.
To see this, observe that there is an action of $\pi_1(\mathscr{M}_{g,2}, x)$
on the set of surjective homomorphisms $\pi_1(C^{\circ},x) \to G$ sending $\phi:
\pi_1(C^{\circ},x) \to G$ to the map $\phi^h(g) := \phi(hgh^{-1})$.
By definition, we have $h \in \Gamma$ if and only if $\gamma^h$ is conjugate to $\gamma$.
In particular, $\Gamma$ contains the stabilizer of $\gamma$ under the action of 
$\pi_1(\mathscr{M}_{g,2}, x)$.
But this stabilizer has finite index in 
$\pi_1(\mathscr{M}_{g,2}, x)$ because 
$G$ is finite and $\pi_1(C^\circ, x)$ is finitely generated. 
Therefore,
there are only
finitely many 
homomorphisms $\pi_1(C^\circ, x) \to G$,
and, in particular, finitely many such surjective homomorphisms.

\end{proof}

\subsubsection{}\label{subsubsection:proof-kodaira-parshin}\begin{proof}[Proof of \autoref{proposition:kodaira-parshin}]
\source{geometric-local-systems-general-curves-isomonodromy}
Let $\widetilde{\Gamma}$ be the kernel of the map $\widetilde{\gamma}$ from
\autoref{lemma:center-free-group}.
The subgroup $\widetilde{\Gamma}$ corresponds to a finite \'etale cover
$\mathscr X^\circ \to \mathscr M_{g,2}$. 
Observe that $\mathscr M_{g,2} \subset \mathscr{M}_{g,1} \times_{\mathscr M_g} \mathscr{M}_{g,1}$
can be viewed as a dense open substack, and
let $\mathscr{X}$ be the normalization
of $\mathscr{M}_{g,1} \times_{\mathscr M_g} \mathscr{M}_{g,1}$ in the function
field of $\mathscr{X}^\circ,$ forming the following cartesian diagram
\begin{equation}
	\label{equation:}
\source{geometric-local-systems-general-curves-isomonodromy}
	\begin{tikzcd} 
		\mathscr X^\circ \ar {r} \ar {d} & \mathscr X \ar {d} \\
		\mathscr M_{g,2} \ar {r} & \mathscr{M}_{g,1} \times_{\mathscr
		M_g} \mathscr{M}_{g,1}.
\end{tikzcd}\end{equation}
Restricting the natural map $\mathscr{X}\to\mathscr{M}_{g,1} \times_{\mathscr M_g} \mathscr{M}_{g,1}$ to a fiber $C$ of the universal curve $\mathscr{M}_{g,1} \times_{\mathscr M_g} \mathscr{M}_{g,1} \to \mathscr M_{g,1}$ yields a finite disjoint union of $G$-covers of $C$, ramified only over the tautological marked point of $C$.
We then take our desired relative curve $f: \mathscr X \to \mathscr{M}_{g,1} \times_{\mathscr M_g}
\mathscr{M}_{g,1} \to \mathscr M_{g,1}$ as the resulting
composition.
\end{proof}

In order to use \autoref{proposition:kodaira-parshin}, we will need to know
there are groups $G$ satisfying its hypotheses. We now provide such an example.
\begin{example}
\source{geometric-local-systems-general-curves-isomonodromy}
\notready
	\label{example:kodaira-parshin}
\source{geometric-local-systems-general-curves-isomonodromy}
	As a concrete example of a group $G$ to which 
	\autoref{proposition:kodaira-parshin} applies, we can take $G = S_3$ to be the
	symmetric group on three letters and identify $\pi_1(Y^\circ, y)$
with the free group on the generators $a_1, \ldots, a_g, b_1, \ldots, b_g$. The group $\pi_1(Y, y)$
is generated by $a_1, \ldots, a_g, b_1, \ldots, b_g$ with the relation
$\prod_{i=1}^g \left[ a_i, b_i \right]$.
Consider the surjection $\phi: \pi_1(C^\circ, y)\twoheadrightarrow S_3$
sending $a_1\mapsto (12), b_1 \mapsto (13)$ and sending $a_i \mapsto
\mathrm{id}, b_i
\mapsto \mathrm{id}$ for $i > 1$.
The loop around the puncture maps to 
$\phi(\prod_{i=1}^g \left[ a_i, b_i \right]) = (12)(13)(12)^{-1}(13)^{-1} =
(123) \neq \mathrm{id}$.
\end{example}


\subsubsection{}
\label{subsubsection:proof-counterexample}
\source{geometric-local-systems-general-curves-isomonodromy}
\begin{proof}[Proof of \autoref{theorem:counterexample}]
Let $f: \mathscr X \to \mathscr{M}_{g,1}$ denote the map from
\autoref{proposition:kodaira-parshin}. 
Concretely, we can take $G = S_3$ and the map $\phi$ as in 
\autoref{proposition:kodaira-parshin} to be that given in
\autoref{example:kodaira-parshin}.
Define the local system $\mathbb V := R^1 f_* \mathbb C$ on $\mathscr M_{g,1}$,
and define $\mathscr{F}$ to be the vector bundle $\mathbb{V}\otimes
\mathscr{O}$. Note that $\mathscr{F}$ admits a natural (Gauss-Manin) connection $\on{id}\otimes d$. The local system $\mathbb{V}$ evidently underlies a variation of Hodge structure.

 Let $C$ be a
 fiber of the forgetful morphism $\mathscr{M}_{g,1}\to \mathscr{M}_{g}$.
 Let $X := f^{-1}(C) \subset \mathscr X$. We claim that the flat vector bundle $(\mathscr{F}, \nabla)$ satisfies the
conditions of \autoref{theorem:counterexample}, i.e.~it has semisimple monodromy and $\mathscr{F}|_C$ is not semistable.

We first check that $(\mathscr{F}, \nabla)|_C$ has semisimple monodromy. This is
true for any flat vector bundle arising from the Gauss-Manin connection on the
cohomology of a family of smooth proper varieties, by work of Deligne
\cite[Th\'eor\`eme 4.2.6]{deligne:hodge-ii}.

We now check that $(\mathscr{F}, \nabla)|_C$ is not semistable.  
By \cite[Theorem
4]{cataneseD:answer-to-a-question-by-fujita},
if $X \to C$ is not isotrivial,
$f_* \omega_f$ is a destabilizing subsheaf of $\mathscr F$.
It remains to show $X \to C$ is not isotrivial.
The fiber of $f|_X$ over a point $x\in C$ is a finite disjoint union of
finite covers of $C$, branched only over $x$. These fibers must vary in moduli as $x$ varies,
as there are only finitely many non-constant maps between any
two curves over of genus at least $2$, 
by de Franchis' theorem. 
(See \cite{de1913teorema} or \cite[Corollary 3, p.~75]{samuel1966lectures}, for
example.)
% or	\cite[Theorem	1]{kobayashiO:meromorphic-mappings-onto-compact-complex-spaces-of-general-type}.
%\begin{lemma}
%	\label{lemma:}
%	If $Y$ and $C$ are two curves of genus at least $2$, there are only
%	finitely many maps $Y \to C$.
%\end{lemma}
%\begin{proof}
%	First, note that the degree of any map $Y \to C$ is bounded by at most
%$\frac{g(Y)-1}{g(C)-1}$ by
%	Riemann-Hurwitz.
%	Next, observe that the moduli space of maps of fixed degree $d$ from $X$ to $C$ is of finite
%	type. Indeed, by sending a map to its graph, we can identify this moduli
%	space of maps with a Hilbert scheme of subschemes
%	$Y\subset X \times C$ so that the projection of $Y$ onto $X$ is an
%	isomorphism and the projection onto $C$ has degree $d$.
%	Therefore, it is enough to show that the moduli space of maps from $X$
%	to $C$ has trivial tangent space at any such map $[f]$. The tangent
%	space is identified with $H^0(X, f|_X^* T_C)$, which vanishes as $f|_X^* T_C$ has negative degree when the genus
%of $C$ is at least $2$. \end{proof}
\end{proof}
\begin{remark}
\source{geometric-local-systems-general-curves-isomonodromy}
\notready
	\label{remark:alternate-hodge-proof}
\source{geometric-local-systems-general-curves-isomonodromy}
	We can also give a somewhat more involved proof of
	\autoref{theorem:counterexample} using \autoref{corollary:unstable} in
	place of \cite[Theorem
	4]{cataneseD:answer-to-a-question-by-fujita}, as we now explain. This argument inspired the Hodge-theoretic results \autoref{theorem:very-general-VHS} and \autoref{corollary:geometric-local-systems}, proven in \autoref{section:hodge-theoretic-results}.

	With notation as in the proof of 
	\autoref{theorem:counterexample},
	$\mathscr{F}$ has degree $0$ since it admits a flat connection, by
	\autoref{prop:basic-facts}(3).
	Therefore, it suffices to show $\mathscr{F}$ has a subsheaf of positive
degree.
The Hodge filtration exhibits $F^1\mathscr F\simeq (f|_X)_* \omega_{(f|_X)}$ as a subsheaf
of $\mathscr F$,
which is destabilizing by
\autoref{corollary:unstable}
once we verify that
$\delta: F^1\mathscr F \to F^0 \mathscr F/ F^1 \mathscr F \simeq R^1 (f|_X)_* \mathscr{O}_{X}$
is nonzero.

We now check $\delta$ is nonzero.
Locally around a point of $C$, $\delta$ 
can be identified with the derivative of the period map
\cite[Theorem 5.3.4]{carlsonMP:period-mappings}
sending a curve corresponding to a fiber of 
$f|_X : X \to C$
to the corresponding Hodge structure on its first cohomology. To show this derivative is not identically zero it suffices to show that the period map is non-constant.
More concretely, by the Torelli theorem, we only need to check
$f|_X: X \to C$ is not isotrivial.
This follows by de Franchis' theorem, as explained in the proof of
\autoref{theorem:counterexample}.
\end{remark}

\begin{remark}
\source{geometric-local-systems-general-curves-isomonodromy}
\notready
	\label{remark:bhh-error}
\source{geometric-local-systems-general-curves-isomonodromy}
	As noted prior to its statement,
	\autoref{theorem:counterexample} contradicts
\cite[Theorem 1.3]{BHH:logarithmic},
\cite[Theorem 1.3]{BHH:irregular}, and
\cite[Theorem 1.2]{BHH:parabolic}, which claim, for example, that any irreducible flat vector bundle on a smooth proper  curve of genus at least $2$ admits a semistable isomonodromic deformation. We now explain the gaps in the proofs of
those results. The error in \cite[Theorem 1.3]{BHH:logarithmic} occurs in
\cite[Proposition 4.3]{BHH:logarithmic}; the issue is that the map denoted
$f^*\nabla$ in diagram (4.14) does not in general exist. The proof works
correctly if $G=\on{GL}_2$. An identical error occurs in \cite[Proposition
4.4]{BHH:irregular}. A different argument is given in \cite{BHH:parabolic}.
There, the error occurs in the proof of \cite[Proposition 5.1]{BHH:parabolic}, 
in which the large diagram claimed to be commutative does not in general commute.
\end{remark}
